\section*{الفصل \ref{ch:b1:nucleophilic-substitution} --- الاستبدال المحب للنواة}

\begin{solution}{exo:b1:nucleophilic-substitution:1}
\ce{CH3CH2OH + Br-}؛ \ce{CH3OCH2CH3 + I-}؛ \ce{CH3CH2CH2CN + Cl-}؛
\ce{CH3NH3+ + Br-}. الركائز: الهالوجينو ألكانات؛ \omterm{def:b1:electronic-effects:nucleophile}{والمحبات للنواة}:
\ce{OH-}، \ce{CH3CH2O-}، \ce{CN-}، \ce{NH3}؛ \omterm{def:b1:electronic-effects:nucleophile}{والمجموعات المغادرة}: أيونات الهاليد.
\end{solution}

\begin{solution}{exo:b1:nucleophilic-substitution:2}
الرتبة الثانية: تتضاعف ثلاث مرات؛ وتنصيف التركيزين كليهما يقسمها على 4. والرتبة الأولى بالنسبة إلى RY وحده:
لا تتغير؛ تتنصف.
\end{solution}

\begin{solution}{exo:b1:nucleophilic-substitution:3}
SN2 (أولي، \omterm{def:b1:electronic-effects:nucleophile}{محب للنواة} جيد، لابروتوني)؛ SN1 (ثالثي، ماء)؛ SN2
(ميثيلي، \omterm{def:b1:electronic-effects:nucleophile}{محب للنواة} قوي)؛ \omterm{def:b1:nucleophilic-substitution:sn1}{تحلل بالمذيب} SN1 (ثانوي، \omterm{def:b1:electronic-effects:nucleophile}{محب للنواة} ضعيف
متعادل، \omterm{def:b1:intermolecular-forces-solvents:solvent-classes}{مذيب بروتوني})، وربما مع شيء من SN2.
\end{solution}

\begin{solution}{exo:b1:nucleophilic-substitution:4}
يحتل السيانيد الموضع المقابل للبروم وتنقلب المجموعات الثلاث
الأخرى. وفي الناتج تحتل المجموعة CN الرتبة الأولى، كما كان البروم،
والرتب الأخرى لا تتغير: فللترتيب المقلوب الواصف
المعاكس، 2-ميثيل بوتان نتريل-\cip{S}.
\end{solution}

\begin{solution}{exo:b1:nucleophilic-substitution:5}
1. \omterm{def:b1:electronic-effects:cleavage}{الانشطار غير المتجانس}: سهم من الرابطة C--Br إلى Br، يعطي \ce{(CH3)3C+}
و\ce{Br-} (بطيء). 2. \omterm{def:b1:lewis-resonance-vsepr:lewis}{زوج حر} للماء نحو الكربون الكاتيوني:
\ce{(CH3)3C-OH2+}. 3. \omterm{def:b1:lewis-resonance-vsepr:lewis}{زوج حر} لجزيء ماء آخر يأخذ بروتونًا
من المجموعة \ce{OH2+}، ويعود زوج O--H إلى الأكسجين. ولا يحدد السرعة إلا الخطوة
الأولى البطيئة؛ والماء، وهو المذيب، بفائض كبير إلى حد أن
تركيزه لا يتغير.
\end{solution}

\begin{solution}{exo:b1:nucleophilic-substitution:6}
الكربون الحامل لذرة Br أولي، لكن بجواره مجموعة ثالثي البوتيل تملأ
الحيز خلف الرابطة C--Br: \omterm{def:b1:elementary-steps:energy-profile}{فالحالة الانتقالية} SN2 مزدحمة
جدًا. وكانت SN1 ستحتاج إلى \omterm{def:b1:electronic-effects:intermediates}{كاتيون كربوني} أولي، غير ثابت إلى حد كبير.
\end{solution}

\begin{solution}{exo:b1:nucleophilic-substitution:7}
يحل اليوديد محل اليوديد \omterm{def:b1:nucleophilic-substitution:sn2}{بالآلية SN2}، مع انقلاب في كل مرة: فيصير \cip{S}
\cip{R} حتى يتساوى وجودهما، وهذا \omterm{def:b1:stereochemistry-in-depth:enantiomers}{خليط راسيمي}. وكل
استبدال يحوّل جزيئًا \cip{S} واحدًا إلى جزيء \cip{R} واحد: فيهبط الدوران
بما يعادل جزيئين، ولذلك يتناقص الدوران أسرع بمرتين من
استبدال \omterm{def:b1:nucleophilic-substitution:substitution}{الركيزة}.
\end{solution}

\begin{solution}{exo:b1:nucleophilic-substitution:8}
SN2: \babelsublr{\foreignlanguage{arabic}{البروموميثان} $>$ \foreignlanguage{arabic}{1-بروموبروبان} $>$ \foreignlanguage{arabic}{2-بروموبروبان} $>$
\foreignlanguage{arabic}{2-برومو-2-ميثيل البروبان}} (الازدحام). SN1: الترتيب المعاكس (ثبات
\omterm{def:b1:electronic-effects:intermediates}{الكاتيون الكربوني})، والبروموميثان و1-بروموبروبان لا يكادان يتفاعلان.
\end{solution}

\begin{solution}{exo:b1:nucleophilic-substitution:9}
$x_{\text{منقلب}} - x_{\text{محتفظ}} = 0.12$ و$x_{\text{منقلب}} + x_{\text{محتفظ}} = 1$:
\qty{56}{\percent} منقلب، و\qty{44}{\percent} محتفظ بترتيبه. فهو \omterm{def:b1:nucleophilic-substitution:sn1}{تحول راسيمي} SN1
في الغالب، مع فائض طفيف من الانقلاب: فالبروميد المغادر
ما زال يحجب جهته حين يصل الماء.
\end{solution}

\begin{solution}{exo:b1:nucleophilic-substitution:10}
$v = k_1k_2[\mathrm{RY}][\mathrm{Nu}]/(k_{-1}[\mathrm{Y^-}] + k_2[\mathrm{Nu}])$
(\cref{prop:b1:nucleophilic-substitution:sn1-law}). وإضافة \ce{Y-}
تزيد المقام: فيعود مزيد من \omterm{def:b1:electronic-effects:intermediates}{الكاتيونات الكربونية} إلى \omterm{def:b1:nucleophilic-substitution:substitution}{الركيزة}
بدل أن تُؤسر، وتنخفض السرعة. ويكون التأثير مهملًا حين يكون
$k_2[\mathrm{Nu}] \gg k_{-1}[\mathrm{Y^-}]$، كما حين يكون \omterm{def:b1:electronic-effects:nucleophile}{المحب للنواة} هو
المذيب.
\end{solution}

\begin{solution}{exo:b1:nucleophilic-substitution:11}
تركيزان متساويان: $-\mathrm{d}c/\mathrm{d}t = kc^2$، $1/c = 1/C_0 + kt$،
$t_{1/2} = 1/(kC_0) = \qty{5.0e5}{s}$ (نحو ستة أيام). ومع فائض
من الهيدروكسيد بخمسين ضعفًا، يكون $[\ce{OH-}] \approx 50C_0$ ثابتًا: $c = C_0
e^{-k't}$ مع $k' = 50kC_0 = \qty{1.0e-4}{s^{-1}}$، $t_{1/2} = \ln 2/k' =
\qty{6.9e3}{s}$، أي نحو ساعتين.
\end{solution}

\begin{solution}{exo:b1:nucleophilic-substitution:12}
كان على البروميد أن يطرد \ce{OH-}، وهو \omterm{def:b1:predominant-reaction:strong-acid}{قاعدة قوية} \omterm{def:b1:electronic-effects:nucleophile}{ومجموعة مغادرة}
رديئة: فلا تفاعل. وفي HBr المركّز تتبرتن OH إلى
\ce{-OH2+}، ومجموعتها المغادرة الماء. الكحول الثالثي: يغادر الماء،
فيتكوّن \omterm{def:b1:electronic-effects:intermediates}{الكاتيون الكربوني}، الذي يأسره \ce{Br-} (SN1). الكحول الأولي:
يهاجم \ce{Br-} كربون \ce{R-OH2+} من الخلف بينما يغادر الماء
(SN2).
\end{solution}

\begin{solution}{pb:b1:nucleophilic-substitution:1}
\textbf{1.} كل جزيء يتفاعل يعطي أيون \ce{H+} واحدًا وأيون \ce{Cl-} واحدًا؛
والموصلية مجموع الإسهامات الأيونية، فهي متناسبة مع
كمية الأيونات.
\textbf{2.} $[\mathrm{RCl}]/[\mathrm{RCl}]_0 = (\kappa_\infty -
\kappa)/\kappa_\infty$.
\textbf{3.} $\ln[\mathrm{RCl}] = \ln[\mathrm{RCl}]_0 - kt$: ارسم
$\ln(\kappa_\infty - \kappa)$ بدلالة $t$.
\textbf{4.} 0.693، 0.603، 0.513، 0.423، 0.333، 0.153، $-0.117$، $-0.387$.
\textbf{5.} نعم، الميل $-0.0180$ في الدقيقة على طول المنحنى: الرتبة الأولى.
\textbf{6.} أي رتبة بالنسبة إلى الماء: فتركيزه ثابت
(انحطاط الرتبة).
\textbf{7.} $k = 0.0180/60 = \qty{3.0e-4}{s^{-1}}$.
\textbf{8.} $t_{1/2} = \ln 2/k = \qty{2.3e3}{s}$، أي \qty{38.5}{min}.
\textbf{9.} $\ln 10/k = \qty{7.7e3}{s}$، أي \qty{128}{min}.
\textbf{10.} $\ln(\kappa_\infty - \kappa)$: 0.693 عند $t = 0$، و0.109 عند
\qty{10}{min}، وهكذا: الميل $-0.0584$ في الدقيقة، $k = \qty{9.7e-4}{s^{-1}}$.
\textbf{11.} $9.7/3.0 = 3.2$.
\textbf{12.} عند \qty{30}{min}: $2.000(1 - e^{-0.0584 \times 30}) = 1.653$،
والمقيس 1.654.
\textbf{13.} SN1: \omterm{def:b1:nucleophilic-substitution:substitution}{ركيزة} ثالثية، \omterm{def:b1:intermolecular-forces-solvents:solvent-classes}{ومذيب بروتوني}، \omterm{def:b1:electronic-effects:nucleophile}{ومحب للنواة} متعادل.
\textbf{14.} كما في \cref{exo:b1:nucleophilic-substitution:5}، مع Cl.
\textbf{15.} لا يشارك في التأين البطيء إلا \omterm{def:b1:nucleophilic-substitution:substitution}{الركيزة}: $v =
k[\mathrm{RCl}]$.
\textbf{16.} لا يمكن للهيدروكسيد أن يؤثر إلا في خطوة الأسر السريعة، بعد
التأين المحدد للسرعة؛ فالسرعة لا تتعلق به.
\textbf{17.} \omterm{def:b1:electronic-effects:intermediates}{الكاتيون الكربوني} مستوٍ \omterm{def:b1:stereochemistry-in-depth:stereocentre}{ولاكيرالي}؛ ويضاف الماء على الوجهين
بالتساوي: كحول راسيمي.
\textbf{18.} حاجز أول مرتفع (التأين) نحو \omterm{def:b1:electronic-effects:intermediates}{الكاتيون الكربوني} في
بئر ضحلة، ثم حاجز منخفض نحو الأسر: فالخطوة الأولى هي
المحددة للسرعة.
\textbf{19.} كان سيكوّن \omterm{def:b1:electronic-effects:intermediates}{كاتيونًا كربونيًا} أوليًا (غير ثابت إلى حد كبير)، والماء
\omterm{def:b1:electronic-effects:nucleophile}{محب للنواة} أضعف من أن يعطي SN2 سريعًا.
\textbf{20.} $k = A\,e^{-E_a/RT}$.
\textbf{21.} $E_a = R\ln(k_2/k_1)/(1/T_1 - 1/T_2)$.
\textbf{22.} $E_a = 8.314 \times \ln(9.74/3.00)/(1/298.15 - 1/308.15) =
\qty{9.0e4}{J/mol}$.
\textbf{23.} ارتفاع حاجز التأين، وهو \omterm{def:b1:elementary-steps:rds}{الخطوة
المحددة للسرعة}.
\textbf{24.} $k = 3.0 \times 10^{-4} \exp\bigl(\frac{90\,000}{8.314}
(\frac{1}{298.15} - \frac{1}{318.15})\bigr) = \qty{2.9e-3}{s^{-1}}$.
\textbf{25.} \textbf{$E_a = \qty{90}{kJ/mol}$}.
\end{solution}
